\chapter{\MakeUppercase{Methodology and Testing}}

\section{Experimental Methodology}
The evaluation compares Epidemic, Spray and Wait, and Spatio-Semantic routing under three traffic levels. Each protocol is executed for 20 runs per traffic level. The simulator averages the resulting metrics and writes one \ac{csv} file per protocol. This produces a direct comparison because all protocols are evaluated using the same node counts, area size, buffer size, contact range, message \ac{ttl}, and traffic probabilities.

\section{Experimental Controls and Fairness}
The methodology is designed to ensure that the comparison is fair. Every protocol is evaluated under the same disaster area, node composition, mobility model, buffer size, message \ac{ttl}, and traffic probabilities. The proposed router is allowed to use only the local state maintained during the simulation, not any external oracle information. This is important because a protocol should not appear superior merely because it receives stronger assumptions than the baselines.

The test campaign also treats protocol outputs uniformly. Delivery, delay, overhead, and drops are computed from the same event counters regardless of which router is active. Result aggregation is performed after repeated runs so that random variation in node movement and message generation does not dominate the conclusions. The intent is to compare routing strategies, not isolated lucky or unlucky runs.

\section{Traffic Scenarios}
The traffic scenarios are defined by per-node message generation probability:

\begin{table}[h!]
	\centering
	\caption{Traffic scenarios}
	\renewcommand{\arraystretch}{1.3}
	\begin{tabular}{|l|l|l|}
		\hline
		\textbf{Traffic Level} & \textbf{Generation Probability} & \textbf{Interpretation} \\
		\hline
		Low & 3/3600 & Light background communication \\
		\hline
		Medium & 8/3600 & Moderate disaster coordination traffic \\
		\hline
		High & 20/3600 & Congested emergency communication \\
		\hline
	\end{tabular}
	\label{tab:traffic_scenarios}
\end{table}

\section{Baseline Protocol Testing}
Epidemic routing is tested to measure the behavior of unrestricted replication. It forwards every message to every encountered node that does not already have it, provided the receiver buffer has space. This baseline helps show the cost of blind reachability.

Spray and Wait routing is tested to measure copy-limited forwarding. The implementation gives one copy at a time during the spray phase and forwards a single-copy message only when the destination is encountered. This baseline helps show the cost of conservative forwarding.

\section{Proposed Protocol Testing}
Spatio-Semantic routing is tested using the same traffic and mobility conditions. Its main test objective is to determine whether spatial, semantic, and priority-aware decisions improve performance under load. The protocol is expected to be especially strong when traffic is high because congestion exposes whether buffer and relay choices are meaningful.

In addition to aggregate performance, the proposed protocol is interpreted in relation to its design claims. If the critical buffer reserve is useful, critical delivery should degrade more slowly than in the flooding baseline as traffic increases. If semantic and spatial relay selection are useful, average critical delay should improve compared with Spray and Wait, especially once simple waiting becomes costly. These expectations provide a more detailed basis for reading the results than a single overall winner/loser comparison.

\section{Dashboard and Visual Testing}
The project includes a Streamlit dashboard for plotting metric trends, radar comparisons, raw results, and launching the live visual simulation. The live visualizer compares the three protocols side by side with the same initial seed and displays delivery, critical delivery, critical delay, overhead, drops, contacts, node roles, and critical-message carriers.

\begin{figure}[h!]
    \centering
    \fbox{\parbox{0.88\textwidth}{\centering
    Image placeholder: add a screenshot of the Streamlit Metric Trends tab showing delivery ratio, critical delivery ratio, delay, overhead, and buffer drops.}}
    \caption{Dashboard metric trends for the three routing protocols}
    \label{fig:dashboard_trends_placeholder}
\end{figure}

\begin{figure}[h!]
    \centering
    \fbox{\parbox{0.88\textwidth}{\centering
    Image placeholder: add a screenshot of the Performance Radar tab for the high traffic scenario.}}
    \caption{Radar comparison of multi-objective routing performance}
    \label{fig:radar_placeholder}
\end{figure}

\begin{figure}[h!]
    \centering
    \fbox{\parbox{0.88\textwidth}{\centering
    Image placeholder: add a screenshot of the three-panel live simulation with Epidemic, Spray and Wait, and Spatio-Semantic views.}}
    \caption{Live three-panel simulation of competing DTN routing protocols}
    \label{fig:live_sim_placeholder}
\end{figure}

\section{Testing Criteria}
The protocol is considered successful if it satisfies the following criteria:

\begin{itemize}
	\item It should match or exceed Epidemic routing in delivery under light and medium traffic without increasing congestion unnecessarily.
	\item It should significantly outperform Epidemic routing under high traffic, where flooding is expected to collapse.
	\item It should reduce critical-message delay compared with Spray and Wait, because emergency messages must not wait passively for direct delivery.
	\item It should maintain a clear advantage in critical delivery ratio under congested conditions.
\end{itemize}

\section{Reproducibility and Validity Considerations}
Running 20 simulations per scenario improves robustness by reducing the effect of any single mobility trace or contact pattern. Averaging is especially important in \acp{dtn}, where small changes in early encounters can influence later message spread. By repeating each configuration multiple times, the methodology produces results that are more representative of the protocol behavior than a single run would provide.

There are still validity limits that should be acknowledged. The random waypoint mobility model does not capture all real disaster movement patterns, and the simulated wireless medium does not include interference, channel contention, or energy depletion. These limitations do not invalidate the comparison, but they define its scope. The testing methodology is best interpreted as a controlled evaluation of routing logic under abstract disaster conditions rather than as a direct prediction of field performance.

\section{Result Interpretation Strategy}
The evaluation is interpreted using both absolute values and comparative trends. Absolute values are useful because they show the practical size of delivery, delay, overhead, and drop behavior under each traffic condition. Comparative trends are equally important because they reveal how rapidly each protocol degrades as the network moves from low load to congestion. This thesis relies on both views because a protocol may appear acceptable in one isolated scenario while still behaving poorly as stress increases.

Special attention is given to high-traffic performance because that scenario is closest to the thesis motivation. In disaster communication, the most difficult periods often occur when the network is busiest and infrastructure support is weakest. A protocol that performs well only when load is light is therefore less convincing for this application. The testing strategy intentionally emphasizes whether the proposed design remains useful when the environment becomes challenging.

\section{Role of Visual Validation}
The dashboard and live simulation are not substitutes for quantitative evaluation, but they provide a useful complementary view of protocol behavior. Plots make it easier to identify broad trends across traffic levels, while the live simulation helps explain how those trends emerge from contact opportunities, buffer pressure, and relay movement. This is particularly valuable in \acp{dtn}, where protocol behavior can be strongly shaped by spatial separation and intermittent encounters.

Visual validation is also helpful for detecting unintuitive behavior. For example, a protocol may show acceptable aggregate delivery while still producing visible congestion around certain node classes or regions of the map. By pairing numerical metrics with interactive views, the methodology improves interpretability and supports stronger discussion in the results chapter.
